\documentclass[11pt]{article}
\usepackage[T1]{fontenc}
\usepackage{amsmath,amssymb,amsfonts}
\usepackage{geometry}
\geometry{margin=1in}
\usepackage{parskip}
\usepackage{xurl}
\usepackage[dvipsnames]{xcolor}
\usepackage{hyperref}
\hypersetup{colorlinks=true, linkcolor=blue, citecolor=blue, urlcolor=blue!70!black, pdfborder={0 0 0}}

\newcommand{\eps}{\varepsilon}

\title{Errata to ``The Universality of Thomas L. Etter's Three-Place Identity:\\
A Formal Proof and Its Significance''}
\author{James Bowery}
\date{\today}

\begin{document}
\maketitle

These corrections concern the appendices and the accompanying Lean~4 modules of the above
paper. They fall \emph{entirely} on the present author's reconstructions and formalization;
\emph{none} of them touches Etter's own constructions, which they leave intact and, in
several places, restore to their full strength. Because a list of corrections read in
isolation can do an injustice to the work it corrects, the standing record is stated first.

\section{What Etter established, and what stands}

Etter's thesis --- that \emph{identity}, taken as a relative, three-place predicate, is
expressively adequate for the whole of mathematics --- is correct, and is of first-rank
foundational significance. It inverts the usual order of foundations: not ``things exist,
then we ask whether they are identical,'' but identity (from a viewpoint) as primitive, with
membership, and hence set theory, \emph{derived}. The following are Etter's, are correct, and
are unaffected by anything below.

\begin{itemize}
\item \textbf{The Universality Theorem (main body).} Definitions D1 and D2, Theorems
4.1--4.3, and the round-trip of \S5 are a faithful transcription of Etter's construction in
\emph{Three-place Identity}~[E3]: membership is definable from three-place identity and
\emph{vice versa}, with the two round-tripping. The mechanization discharges this with a
clean axiom profile. This is the paper's central result and it is correct.

\item \textbf{The three-equality expressiveness theorem (RCV\,$\to$\,ZF).} That ZF set theory
is constructively definable from three equality predicates is \emph{proved in full} by Etter
in \emph{The Expressive Power of Equality}~[E1], via the cell construction. This is a
complete theorem, not a conjecture. Its mechanization (\texttt{Cell.lean}) discharges Etter's
Theorems~1 and~2 with a clean \texttt{\#print axioms}.

\item \textbf{The two-equality Stereo Equality theorem.} A genuine research frontier, not an
error. Etter announces it in \emph{Three-place Identity}~[E3, \S3], defers its proof to
Section~3 of \emph{Membership and Identity}~[E2], and sketches the key \emph{link}
construction in [E3, fn.~5]; the M\&I draft was never completed (so noted therein by
R.~Shoup). It is stated precisely, and recorded as an open problem, below.

\item \textbf{``Universality'' as relative interpretability.} The paper's use of
``universal''/``open system'' is Etter's own rigorous notion of relative interpretability,
and is the correct claim --- expressive adequacy, not model-theoretic foundationalism.
\end{itemize}

\emph{None of the corrections that follow qualifies any of these.}

\section{What the corrections change}

Where the paper went wrong was in the present author's \emph{secondary presentation} of
Etter's two- and three-equality constructions, in Appendices~A and~B and their Lean modules.
There the author substituted a \emph{symmetric} intersection formula --- a reconstruction of
his own, which does not appear in Etter's papers --- for Etter's actual, \emph{directed}
constructions, and a fragmentary mechanization for a complete one. The corrections below
remove that authorial error and \emph{restore} Etter's stronger original work.

\subsection*{Erratum 1 --- the author's intersection membership is not a ZF membership}

\paragraph{As published.} Appendix~A defines $y \in' x \iff \neg\bigl(E_1(y, x) \land
E_2(y, x)\bigr)$, and \S B.3 defines $y \in_{\mathrm{RCV}} x \iff \neg\bigl(R(y, x) \land
C(y, x)\bigr)$, each asserted to satisfy the axioms of ZF.

\paragraph{Correction.} These formulas were introduced by the present author. Because
$E_1, E_2$ (resp.\ $R, C$) are equivalence relations, hence symmetric, each defined relation
is symmetric, $y \in' x \iff x \in' y$. ZF membership is not symmetric ($\emptyset \in
\{\emptyset\}$ but $\{\emptyset\} \notin \emptyset$), so these relations cannot be ZF
memberships, and the assertions that they ``satisfy the axioms of ZF'' are withdrawn. The
defect is one of the author's reconstruction; Etter's own constructions (Erratum~2) are
directed, and are unaffected.

\subsection*{Erratum 2 --- Etter's actual RCV\,$\to$\,ZF construction (restored)}

\paragraph{Correction.} The intersection formulas above do not occur in Etter's surviving
papers. Etter's RCV\,$\to$\,ZF membership is the \emph{cell construction} of \emph{The
Expressive Power of Equality}~[E1]: a cell is an ordered triple $\langle x, y, n\rangle$ with
$x \,\eps\, y$ and $n \in \{1, 2\}$, of value $\mathrm{Val}(\langle x, y, n\rangle) = x$ if
$n = 1$ else $y$; $R, C, V$ are the row, column and value equalities on cells, with $V$ taken
as identity; and membership is recovered from $R, C, V$ alone by
\[ \mathrm{First}(c) \iff \forall d\,\exists e\,\bigl(V(e{=}d) \land C(e{=}c)\bigr), \]
\[ c\,E\,d \iff \exists e, f\,\bigl(R(e{=}f) \land \mathrm{First}(e) \land V(c{=}e) \land
\neg\mathrm{First}(f) \land V(d{=}f)\bigr). \]
The argument-asymmetry supplied by $\mathrm{First}/\neg\mathrm{First}$ is exactly the
directedness the author's symmetric formula lacked. Appendix~B should cite [E1] and use this
construction; the corrected appendix does.

\subsection*{Erratum 3 --- status of the appendix results}

\paragraph{Correction.} The three-equality result --- that ZF is constructively definable in
an RCV language --- is not merely asserted: it is proved in full in [E1] by the cell
construction. The only genuinely open item is the \emph{two}-equality Stereo Equality
theorem. Its construction is the \emph{link sketch} of [E3, fn.~5]: an equality \emph{without
transitivity} $[x{=}y]$ (reflexive and symmetric only), a ``meet'' operator
$[x =: y, z, \dots]$ (``$x$ equals each of $y, z, \dots$ and anything unequal to $y$ is
unequal to $x$''), and a link membership
\[ x \in' y \;:=\; \exists\, x',x'',y',y'',q,q',q''.\;
[q{=:}x,q'] \;\&\; [q'{=:}y,q,q''] \;\&\; [x'{=:}x,x''] \;\&\; [x''{=:}x,x'] \;\&\;
[y'{=:}y,y''] \;\&\; [y''{=:}y,y'], \]
with the claim that there is a set model of $[x{=}y]$ in which $x \in' y$ is ZF membership.
The surviving M\&I draft breaks off before proving this, so the honest status is an
\emph{open conjecture at the frontier} --- not an error, and not a claim the paper was
entitled to assert as proved.

\subsection*{Erratum 4 --- the Lean modules}

\paragraph{Correction.} As published, \texttt{Stereo.lean} and \texttt{RCV.lean} formalized
the author's symmetric intersection formula and proved only disconnected fragments (a trivial
irreflexivity, a definitional link lemma, partial uniqueness); they neither related
membership to the link, nor established any ZF axiom, nor stated the symmetry that refutes
the formula, nor contained Etter's cell apparatus. The main-body round-trip module is
unaffected and stands. A faithful mechanization now accompanies the repository~[E6]:
\texttt{Cell.lean} discharges Etter's Theorems~1 and~2 (the cell construction) over Mathlib's
\texttt{ZFSet}; \texttt{Stereo.intersectionMem\_not\_ZF} records that the author's symmetric
intersection candidate cannot be a ZF membership; and the
two-equality result is stated as a precise, falsifiable proposition
(\texttt{etterStereoConjecture}) transcribing the [E3, fn.~5] link construction.

\paragraph{Net effect.} These corrections \emph{raise} the standing of the program. Where the
paper claimed a result, Etter's own work already supplies a complete proof [E1]; where the
paper reached past what is proved, the honest status is an open conjecture at the frontier of
a still-living research program. The errors were the present author's, in a secondary
presentation; Etter's constructions stand, and are stronger than the reconstruction that
briefly stood in for them.

\section*{References}

\begin{description}
\item[{[E1]}] Etter T. \emph{The Expressive Power of Equality}. Boundary Institute; June 20,
2001. (Contains the cell construction and the proof of the expressiveness/RCV theorem.)
\item[{[E2]}] Etter T. \emph{Membership and Identity}. Boundary Institute; 2006. (Incomplete
draft.)
\item[{[E3]}] Etter T. \emph{Three-place Identity (Set Theory and the Identity of
Mathematical Objects)}. Boundary Institute; 2006.
\item[{[E4]}] Etter T. \emph{Notes on Identity and Pairing}. Boundary Institute; 2001.
\item[{[E5]}] Etter T. \emph{Equalities and Quine Identities}, \S\S1--6. Boundary Institute;
2001.
\item[{[E6]}] Bowery J. \emph{RelativeIdentity} (Lean~4 sources; \texttt{Cell.lean},
\texttt{RCV.lean}, \texttt{Stereo.lean}). \url{https://github.com/jabowery/RelativeIdentity}
\end{description}

\end{document}
